\documentclass[tikz,border=10pt]{standalone}
\usepackage{ctex}
\usepackage{tikz}
\usetikzlibrary{arrows.meta, calc, positioning}

\definecolor{darktext}{RGB}{44,62,80}
\definecolor{nodebg}{RGB}{240,244,248}
\definecolor{nodeborder}{RGB}{52,152,219}
\definecolor{accent}{RGB}{231,76,60}
\definecolor{warn}{RGB}{192,57,43}
\definecolor{grayborder}{RGB}{150,150,150}

\begin{document}
\begin{tikzpicture}[font=\large, line cap=round, line join=round,
    box/.style={draw=nodeborder, thick, fill=nodebg, minimum width=1.0cm, minimum height=0.7cm, rounded corners=3pt, font=\ttfamily, text=darktext},
    arrow/.style={-{Stealth[length=2.5mm]}, thick, nodeborder}
]
  % ========== Title ==========
  \node[text=darktext] at (5.5,6.2) {\textbf{D 题：为什么链表模拟可以变成离线并查集？}};

  % ========== Key insights ==========
  \node[anchor=west, text=accent] at (0.5,5.4) {\textbf{关键洞察：只取最后一次 down 关系}};

  \node[anchor=west, font=\small, text=darktext] at (0.5,4.8) {\textbullet\ \textbf{P} 始终是堆顶 \textcolor{accent}{$\Rightarrow$} 一张卡片最多当一次 \textbf{P}};
  \node[anchor=west, font=\small, text=darktext] at (0.5,4.35) {\textbullet\ \textbf{C} 可以多次移动，但\textbf{只有最后一次}决定最终位置};
  \node[anchor=west, font=\small, text=darktext] at (0.5,3.9) {\textbullet\ 我们不需要知道中间过程，只需要最终的静态森林};

  % ========== Extract down relations ==========
  \node[anchor=west, text=warn, font=\small] at (0.5,3.2) {提取最终关系: down[1]=4, down[4]=2};

  % ========== Forest result ==========
  \node[box] (r2) at (2.5,2.2) {2};
  \node[box] (r4) at (1.5,1.1) {4};
  \node[box] (r1) at (3.5,1.1) {1};
  \node[box] (r3) at (5.5,2.2) {3};
  \node[box] (r5) at (7.5,2.2) {5};

  \draw[arrow] (r4) -- (r2);
  \draw[arrow] (r1) -- (r2);

  \node[text=accent, font=\small] at (7.0,0.3) {并查集合并后，根的大小 = 堆的卡片数};

  % ========== Bottom summary ==========
  \draw[thick, grayborder] (0.2,-0.4) rectangle (10.8,-0.9);
  \node[font=\small, text=darktext] at (5.5,-0.65) {\textbf{本质：最终状态只取决于每个节点最后一次作为 \textbf{C} 被移向哪里}};

\end{tikzpicture}
\end{document}
